\chapter{Injections}

\section*{Overview}
An injection is the administration of a medication or fluid into the body using a needle and syringe. The most common routes are intramuscular (IM) and subcutaneous (SC) injections, which deliver drugs for local or systemic effect.
\section*{Alternative Names}
IM injection; SC injection; subcutaneous injection; intramuscular injection; parenteral drug administration
\section*{Principle}
A drug is deposited into skeletal muscle or subcutaneous fat, where it is absorbed into the circulation. Intramuscular injection provides faster absorption than subcutaneous delivery because of the rich muscle blood supply, while subcutaneous injection is used for drugs requiring slow, sustained absorption.
\section*{History}
The hypodermic syringe was developed in the 19th century, and by the 20th century injection techniques became standard for vaccination, insulin, analgesics, and antibiotics, replacing many older oral and topical routes.
\section*{Why It Is Done}
Ordered to administer vaccines, insulin, anticoagulants, antibiotics, analgesics, hormones, biologics, and other medications when the oral route is unavailable, unreliable, or undesirable, or when rapid or controlled absorption is required.
\section*{Is This a Primary or Secondary Test or Procedure}
A primary therapeutic procedure when the sole purpose is drug administration; it is a delivery technique rather than a diagnostic or curative intervention.
\section*{How It Is Performed}
The site is selected and cleaned. For IM injection the needle is inserted at 90 degrees into a large muscle such as the deltoid, vastus lateralis, or ventrogluteal; for SC injection the needle enters the fatty layer at 45-90 degrees after a skin pinch. After aspirating to confirm the needle is not in a vessel, the medication is injected slowly and the needle is withdrawn with pressure applied.
\section*{Preparation}
Informed consent, verification of the drug, dose, and route, and selection of the appropriate needle length and gauge. The site is assessed for suitability, and the patient's allergy and medication history is reviewed.
\section*{Contraindications and Precautions}
Avoid injecting into an inflamed, infected, bruised, or previously irradiated area, and avoid sites near major nerves or vessels. Use caution in thin or emaciated patients, and aspirate before injection to prevent intravascular delivery.
\section*{Risks}
Pain, bleeding, bruising, infection, abscess, nerve injury, and unintended intravascular injection. Systemic risks include anaphylaxis, adverse drug reactions, and, with insulin, local lipohypertrophy or lipoatrophy.
\section*{Aftercare and Recovery}
The site is observed briefly for immediate reactions, and the patient is monitored for the expected drug effect and any delayed side effects. Injection sites are rotated for repeated administration.
\section*{Outcome and Follow-up}
The response is evaluated by the clinical effect of the medication, such as analgesia, glucose control, or immune response. Failure of effect may reflect incorrect site, dose, or technique.
\section*{Expected Results}
Correct drug delivery with the expected therapeutic response and no significant local or systemic adverse reaction.
\section*{Follow-up}
Patients on repeated injections are taught proper technique and site rotation. Drug levels may be monitored for medications with narrow therapeutic windows, and injection sites are inspected for complications.
\section*{When to Seek Medical Care}
Follow the procedure team's specific instructions. Seek urgent medical assessment for severe or worsening pain, uncontrolled bleeding, fever or signs of infection, trouble breathing, chest pain, fainting, new weakness or confusion, inability to urinate, or any rapidly worsening symptom. The appropriate warning signs depend on the procedure and the patient's medical condition.

\section*{References}
\begin{enumerate}
  \item MedlinePlus. Giving an IM (intramuscular) injection. U.S. National Library of Medicine; 2025.
  \item Current Medical Diagnosis and Treatment. McGraw-Hill; 2025.
\end{enumerate}

\clearpage
